\subsubsection{Near Threshold Kaluza-Klein Graviton Decays and the Dark Dimension --- arXiv:2610.01825}

\textbf{Authors:} Kevin Langhoff

\paragraph{主な主張}
near-threshold の KK graviton tower 内崩壊は pure $d$-wave で、従来の Dark Dimension 文献で仮定されていた $\Gamma\propto p$ ではなく $\Gamma\propto p^5$ と強く抑制されるため、flat dark dimension では KK graviton dark matter を排除し、不可避な freeze-in abundance と MeV gamma-ray 制限を合わせると余剰次元半径を distance conjecture の素朴な期待より約3桁小さくする必要がある \cite{Langhoff:2026NearThreshold}。

\paragraph{新規性}
near-threshold suppression を単なる位相空間効果ではなく leading amplitude が pure $d$-wave になる partial-wave structure として明示し、従来の $s$-wave 見積もりを修正するとともに、その抑制を reheating と gamma-ray cosmology の制限へ直接結びつけた。

\paragraph{位置付け}
同時期の KK graviton cascade 再評価 \cite{Lee:2026KKCascade} と同じ $p^5$ suppression を支持しつつ、partial-wave の物理的解釈と flat Dark Dimension に対する宇宙論的帰結を前面に出した相補的な結果であり、isometry breaking による rapid cascade を前提とする模型を再検討する基準になる。
